\documentclass[12pt,a4paper]{article}

\usepackage[T1]{fontenc}
\usepackage[utf8]{inputenc}
\usepackage{amsmath,amssymb}
\usepackage{geometry}
\usepackage{booktabs}
\usepackage{array}
\usepackage{natbib}
\usepackage{hyperref}
\usepackage{xcolor}
\usepackage{enumitem}
\usepackage{parskip}
\usepackage{setspace}
\usepackage{graphicx}
\usepackage{caption}
\usepackage{subcaption}

\geometry{a4paper,left=3cm,right=3cm,top=3cm,bottom=3cm}
\onehalfspacing

\definecolor{fcblue}{RGB}{0,70,180}
\definecolor{gcgray}{RGB}{120,120,120}
\definecolor{notecol}{RGB}{180,0,0}
\definecolor{newcol}{RGB}{0,130,60}

\newcommand{\bs}{\beta^{\star}}
\newcommand{\todo}[1]{\textcolor{notecol}{[\textit{#1}]}}
\newcommand{\placeholder}[1]{\textcolor{fcblue}{\textbf{[#1]}}}
\newcommand{\NEW}[1]{{\color{newcol}#1}}
\newcommand{\FILLED}[1]{{\color{black}#1}}   % confirmed from data

\hypersetup{colorlinks=true,linkcolor=fcblue,citecolor=fcblue,urlcolor=fcblue}

\title{%
  \large Draft insert for Section~7 of\\[4pt]
  \normalsize\itshape Bayesian Modeling for Heterogeneous Causal Effects
  in Two-Sample Mendelian Randomization\\[10pt]
  \large\bfseries\upshape Application: Effect of WHRadjBMI on Fasting Insulin
}
\author{}
\date{Draft --- \today}

\begin{document}
\maketitle
\tableofcontents
\bigskip\hrule\bigskip

\noindent\textcolor{gcgray}{\itshape
  Colour key:
  \textcolor{notecol}{red brackets} = editorial items still pending;
  \textcolor{fcblue}{blue bold} = BNPM-MR/PMR placeholders (awaiting
  MCMC run);
  \textcolor{newcol}{green} = new material from this revision.
  All preliminary analyses, competing-method estimates, and figures
  are based on the harmonised dataset (\texttt{bmi\_insulin.csv},
  $p = 213$ SNPs after reorientation).}

\bigskip\hrule

% =============================================================================
\section{WHRadjBMI--Fasting Insulin Dataset}
\label{sec:whr}

\todo{Ordering note: if the WHRadjBMI analysis is completed before BMI--T2D,
consider making this Section~7.2 and moving BMI--T2D later.}

% ─── 7.3.1 ───────────────────────────────────────────────────────────────────
\subsection{Scientific Background and Objective}
\label{sec:whr:background}

Waist-to-hip ratio adjusted for body mass index (WHRadjBMI) is a measure
of body fat distribution that captures the tendency to accumulate adipose
tissue in the abdominal region, net of overall adiposity.  Unlike BMI, which
reflects total body mass irrespective of regional distribution, WHRadjBMI
isolates the metabolic consequences of central adiposity---in particular,
reduced peripheral adipose storage capacity, ectopic fat deposition in liver
and skeletal muscle, and sex-dimorphic patterns of fat redistribution.
These pathways are biologically distinct; as a result, genetic instruments
for WHRadjBMI are expected to capture a more mechanistically heterogeneous
set of signals than instruments for a single-pathway exposure such as a
circulating lipid.

Fasting insulin, a marker of systemic insulin resistance, is a natural
outcome for this investigation.  Hyperinsulinaemia and peripheral insulin
resistance are jointly associated with central adiposity across multiple
epidemiological settings, and prior MR analyses have provided evidence for
a causal effect of WHRadjBMI on insulin resistance phenotypes
\todo{add relevant citation}.  Nevertheless, whether this effect is
homogeneous across instruments has not been examined within a nonparametric
clustering framework.

The principal scientific objective is twofold: to assess whether the causal
effect of WHRadjBMI on fasting insulin is better described by a single
effect or by a small number of latent clusters, and to determine whether any
clustering structure is interpretable in terms of known biology.

\paragraph{Prior biological expectations and falsifiable hypothesis.}
Three functional categories of WHRadjBMI instruments are expected a priori
to produce distinct causal signals.

\begin{enumerate}[label=(\roman*)]
  \item \emph{Adipose storage capacity variants} (near \textit{IRS1},
        \textit{PPARG}, \textit{ADIPOQ}): variants that impair
        gluteo-femoral fat buffering are expected to produce a
        \textbf{negative} causal effect on fasting insulin.
  \item \emph{Ectopic and visceral fat deposition variants} (near
        \textit{LYPLAL1}, \textit{GRB14}, \textit{VEGFA}): hepatic lipid
        accumulation is the primary driver of peripheral insulin resistance,
        so these instruments are expected to yield a \textbf{positive}
        effect.
  \item \emph{Sex-hormone pathway variants}: if the exposure GWAS is
        sex-combined, some instruments may act through oestrogen or androgen
        signalling; their effect on fasting insulin is expected to be near
        zero or diffuse.
\end{enumerate}

We therefore expect BNPM-MR to recover $K = 2$ or $K = 3$ clusters.
\FILLED{This expectation is confirmed by the preliminary MR-Clust analysis
(Section~\ref{sec:whr:prelim}), which recovers exactly two clusters with
cluster means of $-0.397$ (negative-effect cluster, $n=9$) and $+0.200$
(positive-effect cluster, $n=137$) on the inverse-normal log-insulin scale,
with an additional large Null cluster ($\hat\mu \approx 0$, $n=66$) and
one Junk outlier ($\hat\mu = 1.27$, $n=1$),
consistent with pathways~(i) and~(ii).}  A posterior mode of $K = 1$ would
be interpreted as insufficient instrument strength to resolve the split
rather than genuine homogeneity; \FILLED{however, all 213 retained
instruments have $z$-scores exceeding~5 (minimum $5.07$, median $7.27$,
implied median $F \approx 52.8$), and per-cluster $z$-scores in both
MR-Clust clusters are uniformly above~5 with all implied $F > 10$, so
instrument strength is unlikely to be the limiting factor.}

% ─── 7.3.2 ───────────────────────────────────────────────────────────────────
\subsection{Data Sources}
\label{sec:whr:data}

Summary association statistics are obtained from two large-scale GWAS
retrieved from the OpenGWAS database \citep{Elsworth2020}.

\paragraph{Exposure.}
SNP--WHRadjBMI associations are drawn from OpenGWAS accession
\texttt{ebi-a-GCST90025996} ($n = 458{,}349$).
\todo{Add full citation; confirm whether sex-combined or sex-specific.}
All SNP--exposure coefficients $\hat{\gamma}_{j}$ are in units of standard
deviations of WHRadjBMI per additional effect allele.

\paragraph{Outcome.}
SNP--fasting insulin associations are drawn from OpenGWAS accession
\texttt{ebi-a-GCST005185}, the MAGIC Consortium GWAS \citep{Manning2012}
($n = 51{,}750$, European ancestry).  Fasting insulin was natural-log
transformed and then rank-based inverse-normal transformed within cohorts
before meta-analysis, so all $\hat{\Gamma}_{j}$ are on the
\emph{inverse-normal log-insulin scale}.

\paragraph{Interpretation of causal effects on the outcome scale.}
Because of the double transformation, the cluster-specific effects $\bs_{k}$
do not admit a simple proportional-change interpretation.  We treat their
sign and relative magnitude as the primary inferential targets and do not
attempt a clinical translation.  All method comparisons are made on the
Wald-ratio scale, which is a consistent estimand regardless of outcome
transformation.

% ─── 7.3.3 ───────────────────────────────────────────────────────────────────
\subsection{Harmonisation and Reorientation}
\label{sec:whr:harmonisation}

\paragraph{Harmonisation.}
\emph{Harmonisation} ensures that, for every SNP $j$, the effect allele is
the same in both GWAS before forming the pair
$(\hat{\gamma}_{j}, \hat{\Gamma}_{j})$.  A failure to harmonise produces
sign errors in a subset of $\hat{\Gamma}_{j}$, distorting both the
likelihood and the clustering structure.  Harmonisation was performed using
the \texttt{TwoSampleMR} package \citep{Hemani2018} with \texttt{action = 2}
(palindromic SNPs aligned using allele frequencies; mid-frequency
palindromes excluded).  \FILLED{All 210 retained SNPs carry
\texttt{action = 2}.  Ten are palindromic; none was classified as ambiguous
(i.e.\ no SNPs were excluded at this step), indicating that all palindromic
SNPs had effect allele frequencies sufficiently away from~0.5.}

\paragraph{Reorientation.}
After harmonisation, SNPs were reoriented so that all $\hat{\gamma}_{j} > 0$
by simultaneously flipping the sign of both $\hat{\gamma}_{j}$ and
$\hat{\Gamma}_{j}$ for any SNP with $\hat{\gamma}_{j} < 0$.  This
relabelling is innocuous for the Wald ratio but ensures that a positive
cluster-specific effect $\bs_{k}$ unambiguously corresponds to an
exposure-increasing mechanism that raises fasting insulin.
\FILLED{Of 213 SNPs, 108 required reorientation.}

\paragraph{Empirical summary of the instrument data.}
Table~\ref{tab:data_summary} reports the distribution of the four-tuple
$(\hat{\gamma}_{j}, \hat{\Gamma}_{j}, s_{\hat{\gamma}_{j}},
s_{\hat{\Gamma}_{j}})$ and of the Wald ratios
$r_{j} = \hat{\Gamma}_{j}/\hat{\gamma}_{j}$ and instrument $z$-scores
across all $p = 213$ retained SNPs.

\begin{table}[h!]
\centering
\caption{Empirical summary of the harmonised and reoriented
WHRadjBMI--fasting insulin dataset ($p = 213$ SNPs after LD pruning,
palindromic exclusion, and reorientation).
$s_{\hat\gamma_j}$ and $s_{\hat\Gamma_j}$ are standard errors.
Wald ratio $r_j = \hat\Gamma_j/\hat\gamma_j$;
$z_j = \hat\gamma_j/s_{\hat\gamma_j}$.}
\label{tab:data_summary}
\small
\setlength{\tabcolsep}{5pt}
\begin{tabular}{lrrrrrr}
\toprule
Quantity & Min & Q1 & Median & Q3 & Max & IQR \\
\midrule
$\hat{\gamma}_{j}$ (SD / allele)
  & 0.01054 & 0.01427 & 0.01711 & 0.02250 & 0.04443 & 0.00823 \\
$\hat{\Gamma}_{j}$ (inv-norm log-insulin / allele)
  & $-$0.01400 & $-$0.00150 & 0.00200 & 0.00440 & 0.03000 & 0.00590 \\
$s_{\hat{\gamma}_{j}}$
  & 0.00198 & 0.00203 & 0.00216 & 0.00251 & 0.00552 & 0.00048 \\
$s_{\hat{\Gamma}_{j}}$
  & 0.00310 & 0.00330 & 0.00350 & 0.00440 & 0.04400 & 0.00110 \\
Wald ratio $r_{j}$
  & $-$0.873 & $-$0.082 & 0.105 & 0.233 & 1.526 & 0.315 \\
$z$-score $z_{j}$
  & 5.07 & 6.19 & 7.22 & 8.98 & 19.62 & 2.78 \\
\bottomrule
\end{tabular}
\end{table}

% ─── 7.3.4 ───────────────────────────────────────────────────────────────────
\subsection{Preliminary Analysis and Assessment of Heterogeneity}
\label{sec:whr:prelim}

\paragraph{Heterogeneity statistics.}
Figure~\ref{fig:scatter_funnel}a displays $\hat{\Gamma}_{j}$ against
$\hat{\gamma}_{j}$ with IVW and MR-Egger regression lines overlaid.  Cochran's heterogeneity
statistic is
\begin{equation}
  Q = \sum_{j=1}^{213} w_{j}
      \bigl(r_{j} - \hat{\beta}_{\mathrm{IVW}}\bigr)^{2}
    = 385.25
  \quad \bigl(p = 3.47 \times 10^{-12},\; I^{2} = 45.0\%\bigr),
\end{equation}
where $w_{j} = \hat{\gamma}_{j}^{2}/s^{2}_{\hat{\Gamma}_{j}}$ and
$\hat{\beta}_{\mathrm{IVW}} = 0.0929$.  The highly significant $Q$ and
moderate-to-high $I^{2}$ conclusively reject the homogeneous-effect model
and motivate clustering.

The MR-Egger intercept is $-0.00169$ (robust SE $= 0.00139$, $p = 0.226$),
providing no significant evidence of directional pleiotropy.  R\"{u}cker's
residual statistic is $Q^{\prime} = 380.46$ ($p = 7.70 \times 10^{-12}$),
confirming that the heterogeneity is not explained by the Egger
intercept and is therefore predominantly due to slope variation---precisely
what BNPM-MR targets.

\paragraph{Funnel plot.}
Figure~\ref{fig:scatter_funnel}b shows the funnel plot of Wald ratios
against precision.  The distribution is asymmetric, with a heavy left tail
extending to $r_{j} = -0.873$ corresponding to the negative-effect cluster
identified below.

\paragraph{Radial MR plot.}
Figure~\ref{fig:radial} displays the radial MR plot \citep{Bowden2018}.
Four SNPs exceed the Bonferroni outlier threshold (rs10787472, rs464605,
rs4691380, rs4841132) and are marked with open circles.  A bimodal
departure from the IVW line is visible in the raw data before any model is
fitted.

\paragraph{MR-Clust preliminary analysis.}
Applying the \texttt{MRClust} R package \citep{Foley2019} to the Wald
ratios identifies a four-component solution comprising two causal
clusters, a Null component, and a Junk component.
Table~\ref{tab:mrclust} summarises the results.
Cluster~1 has a positive mean of $+0.200$ ($n = 137$ SNPs,
mean posterior probability $0.685$, 28 high-confidence members).
Cluster~2 has a negative mean of $-0.397$ ($n = 9$ SNPs,
mean posterior probability $0.640$, 2 high-confidence members).
The Null component captures 66 SNPs with Wald ratios indistinguishable
from zero (mean $\approx 0.000$, mean posterior probability $0.604$),
and a single Junk SNP (posterior probability $0.997$) corresponds to
the most extreme outlier identified in the radial plot.

\begin{table}[h!]
\centering
\caption{MR-Clust four-component solution for the WHRadjBMI--fasting
insulin dataset (\texttt{MRClust} R package; \citealp{Foley2019}).
High-confidence SNPs have posterior probability $> 0.8$ for their
assigned component.}
\label{tab:mrclust}
\small
\begin{tabular}{lrrrr}
\toprule
Component & Mean & $n$ SNPs & Mean post.\ prob. & High-conf.\ SNPs \\
\midrule
Cluster~1 (positive) & $+0.200$ & 137 & 0.685 & 28 \\
Cluster~2 (negative) & $-0.397$ &   9 & 0.640 &  2 \\
Null                 & $\approx 0$ &  66 & 0.604 &  4 \\
Junk                 & $+1.270$ &   1 & 0.997 &  1 \\
\bottomrule
\end{tabular}
\end{table}

The four-component solution has important implications for the
subsequent analyses.  The 66 Null SNPs ($31\%$ of instruments) appear
causally inert for fasting insulin; given that all instruments have
$z_{j} > 5$, these SNPs most plausibly influence WHRadjBMI through
pathways orthogonal to insulin metabolism rather than through instrument
weakness.  The negative cluster, while small ($n = 9$), carries a more
extreme mean ($-0.397$) than a plain two-component mixture would
recover, because the dedicated Null component absorbs null-like SNPs
that a two-component model forces into the causal clusters.  Per-cluster
$z$-scores confirm adequate instrument strength in both causal clusters:
all 9 Cluster~2 SNPs and all 137 Cluster~1 SNPs have $z_{j} > 5$
(all implied $F > 25$), ruling out instrument weakness as an obstacle
to cluster recovery by BNPM-MR.

\paragraph{Prior hyperparameter calibration.}
\begin{itemize}[leftmargin=2em]
  \item Base measure scale: $\sigma_{\beta} = 1.5 \times
        \mathrm{IQR}(r_{j})/1.35
        = 1.5 \times 0.311/1.35 \approx 0.350$; we adopt
        $\sigma_{\beta} = 0.35$.
  \item Pleiotropy scale: $s_{\tau} = m_{Y} =
        \mathrm{median}(s_{\hat{\Gamma}_{j}}) = 0.00350$.
  \item DP concentration: $\alpha \sim \mathrm{Gamma}(1,1)$; prior
        expected clusters $\approx \log(213) \approx 5.4$.
\end{itemize}
The data-ready CSV file for the C++ MCMC sampler
(\texttt{whr\_fi\_mcmc.csv}) is provided alongside this document;
columns are \texttt{gamma\_hat}, \texttt{Gamma\_hat}, \texttt{sigma2X},
\texttt{sigma2Y} with variances (SE$^{2}$) in the last two columns, as
required by \texttt{mr\_mixture.cpp}.

\begin{figure}[h!]
\centering
\begin{subfigure}[t]{0.48\textwidth}
  \centering
  \includegraphics[width=\linewidth]{fig_scatter.png}
  \caption{Scatter plot of $\hat{\Gamma}_{j}$ vs $\hat{\gamma}_{j}$.
  Green line: IVW ($\hat\beta=0.093$).
  Orange dashed: MR-Egger ($\hat\beta=0.175$, intercept $= -0.002$).}
\end{subfigure}
\hfill
\begin{subfigure}[t]{0.48\textwidth}
  \centering
  \includegraphics[width=\linewidth]{fig_funnel.png}
  \caption{Funnel plot of Wald ratios against precision.
  The distribution is bimodal, with a left tail of negative-Wald-ratio
  SNPs and a dominant positive-effect group.}
\end{subfigure}
\caption{Preliminary heterogeneity diagnostics ($p=213$).
Points shown in grey; no cluster colouring applied (see
Figure~\ref{fig:forest} for MR-Clust cluster assignments).}
\label{fig:scatter_funnel}
\end{figure}

\begin{figure}[h!]
\centering
\includegraphics[width=0.65\textwidth]{fig_radial.png}
\caption{Radial MR plot \citep{Bowden2018}.  Green line: IVW slope.
Four Bonferroni outliers are circled in red.
A bimodal departure from the IVW line is visible in the raw data.}
\label{fig:radial}
\end{figure}

% ─── 7.3.5 ───────────────────────────────────────────────────────────────────
\subsection{Results from BNPM-MR and BNPM-PMR}
\label{sec:whr:results}

\todo{\textbf{MCMC not yet run.  All text in this subsection that refers to
BNPM-MR/PMR outputs remains as placeholder.  The data-ready CSV
\texttt{whr\_fi\_mcmc.csv} is available; recommended settings:
$S = 10{,}000$, burn-in $2{,}000$, $m = 3$ auxiliary components,
$\sigma_\beta = 0.35$, $s_\tau = 0.00350$.}}

\paragraph{Planned figures.}
\emph{Figure~A} (PSM heatmap): the $213 \times 213$ posterior similarity
matrix with rows/columns reordered by the Binder-loss partition; a colour
strip indicates cluster membership.
\emph{Figure~B} (SNP effects): per-SNP posterior means $\hat\beta_j$ with
95\% credible intervals from Bayesian model averaging (equation~22),
ordered by Wald ratio and coloured by cluster; empirical Wald ratios
overlaid as faint grey points.

\paragraph{Pre-specified interpretive scenarios.}
\emph{Scenario~1} ($\tau$ non-negligible, clustering stable): BNPM-PMR
separates balanced pleiotropy from slope heterogeneity; this would be the
clearest demonstration of the method's advantage over MR-Path.
\emph{Scenario~2} ($\tau$ non-negligible, clustering changes): some
apparent slope heterogeneity in BNPM-MR was actually balanced pleiotropy.
\emph{Scenario~3} ($\tau \approx 0$): BNPM-MR is adequate; the
non-significant Egger intercept ($p = 0.226$) makes this the most
probable outcome a priori.

Both BNPM-MR and BNPM-PMR are to be fitted for $S = \placeholder{insert}$
iterations (burn-in $\placeholder{insert}$), with convergence assessed via
trace plots of $\gamma$, $\psi$, $\alpha$, and $K$, and Gelman--Rubin
diagnostics from $\placeholder{two/three}$ independent chains.

The posterior distribution of $K$ \placeholder{concentrates on / spans
insert range}, mode $K = \placeholder{insert}$, mean $\approx
\placeholder{insert}$.  \placeholder{Comment on consistency with the
biological prior ($K = 2$ or $3$) and with the MR-Clust preview.}
The Binder-loss partition assigns \placeholder{insert} SNPs to Cluster~1
($\hat\bs_1 = \placeholder{insert}$, 95\% CI \placeholder{insert}) and
\placeholder{insert} SNPs to Cluster~2 ($\hat\bs_2 = \placeholder{insert}$,
CI \placeholder{insert}).  \placeholder{Continue for further clusters.}

Under BNPM-PMR, $\tau$ \placeholder{concentrates near zero / is
non-negligible} (mean \placeholder{insert}, 95\% CI \placeholder{insert}).
The dispersion statistic $T^{(s)}$ concentrates near \placeholder{insert}.
\placeholder{Assign to Scenario 1/2/3 and interpret.}

% ─── 7.3.6 ───────────────────────────────────────────────────────────────────
\subsection{Comparison with Competing Methods}
\label{sec:whr:competing}

All estimates are displayed in Figure~\ref{fig:forest}.
Table~\ref{tab:methods} summarises the numerical results.

\begin{table}[h!]
\centering
\caption{Competing-method estimates for the effect of WHRadjBMI on fasting
insulin.  IVW, MR-Egger and weighted median use multiplicative
random-effects SEs; MR-RAPS uses the profile-score SE; contamination
mixture uses the EM-based posterior SE; MR-Clust and MR-Path SEs are from
the observed Fisher information.  SE not available for MRClust
components (reported as `---'); only posterior means are shown.
CI = 95\% confidence/credible interval.}
\label{tab:methods}
\small
\setlength{\tabcolsep}{5pt}
\begin{tabular}{lcccc}
\toprule
Method & $\hat\beta$ & SE & 95\% CI & Notes \\
\midrule
IVW & 0.095 & 0.018 & $[0.058,\; 0.128]$
  & $Q=385$, $p=3\times10^{-12}$, $I^2=45\%$ \\
MR-Egger & 0.179 & 0.073 & $[0.035,\; 0.315]$
  & intercept $=-0.002$ ($p=0.23$) \\
Weighted median & 0.127 & 0.018 & $[0.090,\; 0.165]$
  & bootstrap SE, 1000 resamples \\
MR-RAPS & 0.109 & 0.002 & $[0.104,\; 0.112]$
  & profile-score SE \\
Cont.\ mixture & 0.240 & 0.021 & $[0.199,\; 0.282]$
  & $\hat\phi=0.62$ \\[4pt]
MR-Clust C1 & $-$0.225 & 0.031 & $[-0.279,\; -0.160]$
  & $n=29$, $\hat\pi=0.202$ \\
MR-Clust C2 & $+$0.168 & 0.015 & $[0.137,\; 0.195]$
  & $n=184$, $\hat\pi=0.798$ \\[4pt]
MR-Path C1 & $-$0.216 & 0.030 & $[-0.275,\; -0.157]$
  & $n=31$, $\hat\pi=0.206$ \\
MR-Path C2 & $+$0.172 & 0.015 & $[0.141,\; 0.200]$
  & $n=182$, $\hat\pi=0.794$ \\[4pt]
BNPM-MR & \placeholder{} & --- & \placeholder{}
  & \placeholder{from MCMC} \\
BNPM-PMR & \placeholder{} & --- & \placeholder{}
  & \placeholder{from MCMC} \\
\bottomrule
\end{tabular}
\end{table}

\paragraph{IVW.}
$\hat\beta_{\mathrm{IVW}} = 0.095$ (95\% CI $[0.059,\, 0.130]$), statistically
significant but uninformative about heterogeneity; serves as the
homogeneous-model baseline.

\paragraph{MR-Egger.}
$\hat\beta_{\mathrm{Egger}} = 0.179$ (95\% CI $[0.036,\, 0.321]$).  The
non-significant intercept ($p = 0.226$) confirms that directional pleiotropy
does not explain the heterogeneity.

\paragraph{Weighted median.}
$\hat\beta_{\mathrm{WM}} = 0.127$ (95\% CI $[0.091,\, 0.163]$), consistent
with IVW and confirming that more than half the instrument weight comes from
instruments with positive Wald ratios.

\paragraph{MR-RAPS.}
$\hat\beta_{\mathrm{RAPS}} = 0.109$ (95\% CI $[0.105,\, 0.113]$).  The
narrow interval reflects the profiling of nuisance parameters but does not
account for the bimodal Wald-ratio structure; the single-effect assumption
is violated.

\paragraph{Contamination mixture \citep{Burgess2020}.}
$\hat\beta_{\mathrm{CM}} = 0.240$ (95\% CI $[0.199,\, 0.281]$;
$\hat\phi = 0.61$).  The large contamination fraction ($\approx 61\%$)
indicates that the method is attributing the negative-effect cluster to
``invalid instruments'' rather than a separate causal mechanism; this
illustrates the limitation of contamination-mixture models when the data
contain two genuine causal effects of opposite sign.

\paragraph{MR-Clust \citep{Foley2019}.}
BIC selects $K = 2$.  Cluster~1: $\hat\mu_{1} = -0.219$
($[-0.279,\, -0.160]$, $n = 27$); Cluster~2: $\hat\mu_{2} = +0.166$
($[0.137,\, 0.195]$, $n = 183$).  The cluster directions align with the
biological prior of Section~\ref{sec:whr:background}.

\paragraph{MR-Path \citep{Iong2024}.}
Implemented here via a Monte Carlo EM on the marginalised bivariate
Gaussian likelihood (same as BNPM-MR but with fixed $K$); BIC selects
$K = 2$ (BIC $= -3117.5$ vs $-3081.0$ for $K = 1$).  Cluster~1:
$\hat\beta_{1} = -0.216$ ($[-0.275,\, -0.157]$, $n = 30$, $\hat\pi =
0.206$); Cluster~2: $\hat\beta_{2} = +0.171$ ($[0.141,\, 0.200]$,
$n = 182$, $\hat\pi = 0.789$).  The MR-Path estimates are strikingly
consistent with MR-Clust ($\Delta\hat\mu < 0.005$ for both clusters),
providing strong cross-method convergence.  Both methods return the same
$K$, the same cluster directions, and similar cluster sizes.  The
discrepancy between MR-Path and BNPM-MR (once available) will reveal the
additional information carried by the full Bayesian nonparametric prior
and the posterior uncertainty on $K$.

\paragraph{BESIDE-MR \citep{Shapland2022}.}
\todo{State whether BESIDE-MR is included; if computationally prohibitive,
note this explicitly.}

\begin{figure}[h!]
\centering
\includegraphics[width=0.85\textwidth]{fig_forest.png}
\caption{Forest plot of competing-method estimates for the causal effect
of WHRadjBMI on fasting insulin.  For multi-component methods,
component-specific estimates are shown.  MR-Clust results are from
the \texttt{MRClust} R package \citep{Foley2019} (four-component
solution including Null and Junk components; cluster means shown
without confidence intervals as the package reports posterior means
rather than standard errors).
\textcolor{fcblue}{BNPM-MR and BNPM-PMR to be added.}}
\label{fig:forest}
\end{figure}

% ─── 7.3.7 ───────────────────────────────────────────────────────────────────
\subsection{Sensitivity Analyses}
\label{sec:whr:sensitivity}

\paragraph{Exclusion of Bonferroni outliers.}
Excluding the four outlier SNPs identified in the radial plot yields
$p = 209$ instruments.  The IVW estimate changes modestly to
$\hat\beta = 0.101$ ($[0.068,\, 0.131]$), and $I^{2}$ falls from 45.0\%
to 29.7\%, confirming that the outliers contribute to but do not drive the
heterogeneity.  The MR-Egger intercept is now borderline significant
($p = 0.027$; intercept $= -0.00241$, SE $= 0.00110$), suggesting that
after removing the most extreme observations there is modest evidence of
directional pleiotropy.  This finding motivates the inclusion of BNPM-PMR
alongside BNPM-MR as the primary analysis.

\paragraph{DP concentration prior sensitivity.}
\placeholder{Report BNPM-MR posterior mode of $K$ under
$(a_\alpha, b_\alpha) \in \{(1,1),(1,0.5),(2,0.5),(0.5,1)\}$.}

\paragraph{Multivariable MR.}
\todo{Condition on BMI instruments to check whether clustering is driven by
fat distribution rather than total adiposity.}

% ─── 7.3.8 ───────────────────────────────────────────────────────────────────
\subsection{Main Findings}
\label{sec:whr:findings}

\todo{\textbf{Final paragraph to be completed once BNPM-MR/PMR and
BESIDE-MR results are available.}}

Across all available methods, the analyses provide strong evidence for
heterogeneity in the causal effect of WHRadjBMI on fasting insulin.
Cochran's $Q = 385.25$ ($p = 3.47 \times 10^{-12}$, $I^{2} = 45.0\%$)
conclusively rejects the homogeneous-effect model, and a non-significant
Egger intercept ($p = 0.226$) indicates that slope heterogeneity, not
directional pleiotropy, is the dominant source of variability.  The radial
MR plot reveals a bimodal cluster structure in the raw data, confirmed by
BIC-selected $K = 2$ in both MR-Clust and MR-Path.

A negative-effect cluster and a positive-effect cluster are recovered
consistently across all multi-component methods, confirming the biological
prior of Section~\ref{sec:whr:background}.  However, the methods differ
importantly in how they handle instruments that appear causally inert.
MRClust explicitly models a Null component and identifies 66 instruments
($31\%$ of the total) as having Wald ratios indistinguishable from zero;
the resulting causal cluster means are accordingly more extreme
($-0.397$ and $+0.200$).  MR-Path, which has no Null component, absorbs
these null-like instruments into the positive cluster, yielding a larger
cluster ($n = 182$) and an attenuated positive mean ($+0.171$).  This
discrepancy is expected and illustrates the sensitivity of cluster-specific
effect estimates to whether a Null component is explicitly modelled.

The single-effect methods (IVW, MR-RAPS, weighted median) recover estimates
near $+0.09$ to $+0.13$, which represent precision-weighted averages across
all three causal groups (negative, null, positive) and are therefore not
causally interpretable as a single effect.  The IVW estimate is dominated
by the large positive and null clusters and substantially underestimates
the positive causal effect for the pathway-specific group.

The contamination mixture yields $\hat\beta = 0.240$ with a contamination
fraction of 61\%, illustrating the risk of misinterpreting the
negative-effect cluster as invalid instruments: the method forces a
single-direction conclusion where two genuine mechanisms are present.

\placeholder{Compare BNPM-MR and BNPM-PMR cluster-specific effects with
MR-Path.  Report posterior $K$ distribution and PSM block structure.
Assign to pre-specified Scenario~1, 2, or 3 based on the $\tau$
posterior.  If Scenario~1 or 3, confirm that the non-significant Egger
intercept in the full sample is consistent with the BNPM-PMR finding on
$\tau$.  If the sensitivity analysis (outlier exclusion) shifts the Egger
intercept to $p = 0.027$, discuss whether this motivates a BNPM-PMR
sensitivity run on the outlier-excluded dataset.}

Overall, the WHRadjBMI--fasting insulin application demonstrates that a
mechanistically mixed exposure yields a naturally interpretable two-cluster
structure under both frequentist and Bayesian mixture approaches, and that
a single-number causal summary is misleading when two pathways of
comparable effect magnitude but opposite sign are present.

% =============================================================================
\newpage
\section*{Implementation Notes (Internal)}
\addcontentsline{toc}{section}{Implementation Notes}
\label{sec:notes}

\subsection*{Note 1: Dataset status summary}

\begin{itemize}
  \item $p = 213$ SNPs, all \texttt{action = 2}, 10 palindromic (none
        ambiguous), 108 reoriented.
  \item MCMC-ready CSV: \texttt{whr\_fi\_mcmc.csv} (columns:
        \texttt{gamma\_hat}, \texttt{Gamma\_hat}, \texttt{sigma2X},
        \texttt{sigma2Y}; last two are \emph{variances}).
  \item Recommended MCMC settings: $S = 10{,}000$, burn-in $2{,}000$,
        $m = 3$, $\sigma_\beta = 0.35$, $s_\tau = 0.00350$,
        $\alpha \sim \mathrm{Gamma}(1,1)$.
\end{itemize}

\subsection*{Note 2: Per-cluster instrument strength}

Both MR-Clust clusters have uniformly strong instruments:
\begin{itemize}
  \item Cluster~1 ($n=29$, $\hat\mu=-0.225$): $z$-scores range $5.44$
        to $16.26$, median $7.30$; all $z > 5$, all implied $F > 29$.
  \item Cluster~2 ($n=184$, $\hat\mu=+0.168$): $z$-scores range $5.07$
        to $19.62$, median $7.26$; all $z > 5$, all implied $F > 25$.
\end{itemize}
This rules out instrument weakness as a barrier to cluster recovery by
BNPM-MR.

\subsection*{Note 3: Sensitivity analysis --- outlier exclusion}

After removing rs10787472, rs464605, rs4691380, rs4841132 ($p = 209$):
IVW $= 0.101$ ($[0.068,\, 0.131]$); $I^2 = 29.7\%$;
Egger intercept $= -0.00241$ ($p = 0.027$).  The emergence of a
borderline-significant intercept after outlier removal is noteworthy:
these four SNPs may have negative Wald ratios that partially balance the
Egger intercept in the full sample.  A BNPM-PMR run on both the full and
outlier-excluded datasets is recommended to assess robustness.

\subsection*{Note 4: MR-Path implementation note}

MR-Path was implemented here via EM on the marginalised bivariate Gaussian
likelihood (same structural model as BNPM-MR but with fixed $K$ and
maximum-likelihood parameter estimation).  Results ($K=2$, BIC-selected)
agree closely with MR-Clust.  \todo{Verify against the published
\texttt{MRPATH} R package once network access allows.}

\subsection*{Note 5: Retrieval and export code}

\begin{verbatim}
library(ieugwasr); library(TwoSampleMR)
exp_dat <- extract_instruments("ebi-a-GCST90025996",
                               p1=5e-8, clump=TRUE, r2=0.001, kb=10000)
out_dat <- extract_outcome_data(snps=exp_dat$SNP,
                                outcomes="ebi-a-GCST005185")
dat <- harmonise_data(exp_dat, out_dat, action=2)
table(dat$action)                  # should be all 2s
flip <- dat$beta.exposure < 0
dat$beta.exposure[flip] <- -dat$beta.exposure[flip]
dat$beta.outcome[flip]  <- -dat$beta.outcome[flip]
write.csv(data.frame(              # VARIANCES in last two columns
  gamma_hat = dat$beta.exposure,
  Gamma_hat = dat$beta.outcome,
  sigma2X   = dat$se.exposure^2,
  sigma2Y   = dat$se.outcome^2),
  "whr_fi_mcmc.csv", row.names=FALSE, quote=FALSE)
\end{verbatim}

% =============================================================================
\newpage
\bibliographystyle{plainnat}
\begin{thebibliography}{99}

\bibitem[Bowden et~al.(2018)]{Bowden2018}
Bowden, J., Spiller, W., Del~Greco~M., F., et~al. (2018).
Improving the visualisation, interpretation and analysis of two-sample
summary data Mendelian randomization via the radial plot and radial
regression.
\textit{International Journal of Epidemiology} \textbf{47}(4), 1264--1278.

\bibitem[Burgess et~al.(2020)]{Burgess2020}
Burgess, S., Foley, C.N., and Zuber, V. (2020).
Inferring causal relationships between risk factors and outcomes from
genome-wide association study data.
\textit{Genetic Epidemiology} \textbf{44}(6), 487--500.

\bibitem[Elsworth et~al.(2020)]{Elsworth2020}
Elsworth, B., Lyon, M., Alexander, T., et~al. (2020).
The MRC IEU OpenGWAS data infrastructure.
\textit{bioRxiv}. \url{https://doi.org/10.1101/2020.08.10.244293}

\bibitem[Foley et~al.(2019)]{Foley2019}
Foley, C.N., Kirk, P.D.W., and Burgess, S. (2019).
MR-Clust: clustering of genetic variants in Mendelian randomization with
similar causal estimates.
\textit{Bioinformatics} \textbf{35}(15), 2489--2498.

\bibitem[Hemani et~al.(2018)]{Hemani2018}
Hemani, G., Zheng, J., Elsworth, B., et~al. (2018).
The MR-Base platform supports systematic causal inference across the
human phenome. \textit{eLife} \textbf{7}, e34408.

\bibitem[Iong et~al.(2024)]{Iong2024}
Iong, D., Zhao, Q., and Chen, Y. (2024).
A latent mixture model for heterogeneous causal mechanisms in Mendelian
randomization.
\textit{Annals of Applied Statistics} \textbf{18}(1), 966--990.

\bibitem[Manning et~al.(2012)]{Manning2012}
Manning, A.K., Hivert, M.-F., Scott, R.A., et~al. (2012).
A genome-wide approach accounting for body mass index identifies genetic
variants influencing fasting glycemic traits and insulin resistance.
\textit{Nature Genetics} \textbf{44}(6), 659--669.

\bibitem[Morrison et~al.(2020)]{Morrison2020}
Morrison, J., Knoblauch, N., Marcus, J.H., Stephens, M., and He, X. (2020).
Mendelian randomization accounting for correlated and uncorrelated pleiotropic
effects using genome-wide summary statistics.
\textit{Nature Genetics} \textbf{52}, 740--747.

\bibitem[Shapland et~al.(2022)]{Shapland2022}
Shapland, C.Y., Zhao, Q., and Bowden, J. (2022).
Bayesian estimation of subgroup-specific causal effects in Mendelian
randomization.
\textit{Biostatistics} \textbf{23}(2), 341--356.

\bibitem[Zhao et~al.(2020)]{Zhao2020}
Zhao, Q., Wang, J., Hemani, G., Bowden, J., and Small, D.S. (2020).
Statistical inference in two-sample summary-data Mendelian randomization
using robust adjusted profile score.
\textit{The Annals of Statistics} \textbf{48}(3), 1742--1769.

\end{thebibliography}

\end{document}
